\section{Pairwise-Independent Hash (SODA Baseline)}
\label{sec:soda}

\subsection{The Hash Function}
\label{subsec:the-hash-function}

The paper's original locality-preserving hash~\citep[§3.1]{goswami2015approximate}:
\begin{equation}
    h(x) = \bigl(u(\lfloor x/2^K \rfloor) + x\bigr) \bmod 2^K
    \label{eq:soda-hash}
\end{equation}
where $u: [0, U/2^K) \to [0, 2^K)$ is drawn from a pairwise-independent family, stored as two
64-bit coefficients $(a, b)$.

\begin{figure}[h]
    \centering
    \includegraphics[width=0.85\textwidth]{soda_hash_mechanism}
    \caption{SODA hash mechanism: per-block rotation maps keys to $[0, 2^K)$.
    A query range spanning two blocks produces at most two intervals in
    the mapped space.}
    \label{fig:soda-mechanism}
\end{figure}

Properties of the hash:
\begin{itemize}
    \item \textbf{Pairwise-independent:} $\Pr[h(x_1) = h(x_2)] \leq 1/2^K$ regardless of
    key structure.
    This holds for sequential, clustered, or adversarial data.
    \item \textbf{Locality-preserving:} the hash applies a per-block rotation, so
    consecutive keys within the same block map to consecutive positions in $[0, 2^K)$.
    When $\mathcal{L} \leq 2^K$, a query range $[a, b]$ spans at most 2 blocks,
    producing at most 2 intervals in $[0, 2^K)$, each checked with one
    \textsc{IsEmpty} call to ERE.
    (The choice $K = \lceil\log_2(n\mathcal{L}/\varepsilon)\rceil \geq \log_2 \mathcal{L}$
    guarantees this.)
    \item \textbf{Compact:} stores only two 64-bit coefficients $(a, b)$ ---
    $O(1)$ bits, independent of $n$.
\end{itemize}

\subsection{ARE Construction With SODA Hash}

This is a direct implementation of the construction from~\citet[§3]{goswami2015approximate}.
The hash design, its proofs of locality-preservation and pairwise independence, and the
resulting space/FPR bounds are all from the original paper --- our contribution is the
practical implementation.

Combining this hash with ERE (Chapter~\ref{chap:ere}) yields the baseline ARE filter:
\begin{enumerate}
    \item Divide $U$ into blocks of size $2^K$. Let $\text{blockIdx}(x) = \lfloor x / 2^K \rfloor$
    be the block index of key~$x$.
    \item For each block, compute pairwise-independent shift: top $K$ bits of
    $(a \cdot \text{blockIdx}(x) + b)$.
    \item Hash each key via Equation~\eqref{eq:soda-hash}.
    \item Sort, deduplicate, build ERE over $[0, 2^K)$.
    \item Query: hash both endpoints; rotation may split the range into two
    intervals --- check both in ERE.
\end{enumerate}

\begin{center}
\colorbox{orange}{\parbox{0.9\textwidth}{\textbf{TODO:} maybe we do not need to explain how to construct ARE here --- it looks like duplication. Maybe just write: \textit{Combining this hash with ERE (Chapter~\ref{chap:ere}) yields the baseline ARE filter: hash all keys via Eq.~\eqref{eq:soda-hash}, sort and deduplicate, build ERE over $[0, 2^K)$.}}}
\end{center}

\paragraph{Filter properties.}
\begin{itemize}
    \item \textbf{FPR $\leq \varepsilon$ for any data distribution} --- follows from the
    pairwise independence of the hash.
    \item \textbf{No false negatives} --- all images of the query range are checked in ERE.
    \item \textbf{Space:} $K = \lceil \log_2(n\mathcal{L}/\varepsilon) \rceil$, giving
    BPK~$= \log_2(\mathcal{L}/\varepsilon) + O(1)$ --- matching the information-theoretic
    lower bound.
    Including ERE metadata (succinct bit vector $D$), the total cost is
    $\log_2(\mathcal{L}/\varepsilon) + {\sim}2$ bits per key.
\end{itemize}

\subsection{Empirical Validation}
\label{subsec:empirical-validation}

On uniform data ($n = 2^{20}$, $\mathcal{L} = 128$), the measured FPR-vs-BPK curve
tracks the theoretical bound with a gap of ${\sim}2$~BPK (ERE overhead,
Figure~\ref{fig:soda-uniform}).

\begin{figure}[h]
    \centering
    \includegraphics[width=0.85\textwidth]{soda_tradeoff_uniform}
    \caption{FPR vs BPK on uniform keys ($n = 2^{20}$, $\mathcal{L} = 128$). SODA tracks
    the theoretical bound with ${\sim}2$~BPK overhead from ERE metadata.}
    \label{fig:soda-uniform}
\end{figure}
